\chapter{The theorem and its scope}\label{ch:introduction}

\section{The problem}
The volume conjecture compares two objects with very different definitions: a finite root-of-unity knot invariant and the complete hyperbolic metric of a knot complement. Kashaev's finite invariant was introduced in \cite{Kashaev1995}; the asymptotic volume proposal appeared in \cite{Kashaev1997}. The relation with the root-of-unity specialization of the colored Jones polynomial is part of the historical setting \cite{MurakamiMurakami2001}. Here the invariant is fixed by explicit matrices and a specified enhancement. The comparison with that conventional normalization is proved by finite identities, rather than used to replace the source by another state integral.

A long-twist theorem has two limits that must be separated. First choose a diagram and a tuple of sufficiently large integer twists. These integers are then held fixed while the level $N$ tends to infinity. The proof must therefore control the whole finite invariant at every large level, not just a favorable color, a positive coefficient, or a formal solution of a recurrence. In particular, the oscillations responsible for the smaller volume after filling cannot be removed by taking absolute values before the filling transform.

This manuscript treats a compositional family. The local objects have four ordered boundary points, so products and turns can be compiled without choosing a new global decomposition at each stage. The main difficulty is that the exact mixed-color leaf contains a second, signed part, called its \emph{alias}. The compositional estimate controls that part when one leaf is attached to a completed child. It is a global supporting-plane inequality on the actual multiplication domain. Its proof uses a finite collection of rational inequalities with rigorous transcendental enclosures.

\section{Ordered tangles and completed expressions}
An elementary augmented tangle consists of two main strands encircled by one auxiliary crossing circle. Write $L$ for that four-port box and $h=\Turn L$ for its ordered quarter-turn. Product means the literal vertical stack, with the two intermediate strands joined in their recorded order. A turn bends the original boundary strands using the original cup and cap. Neither operation permits an unrecorded permutation of ports.

\begin{definition}[Completed grammar]\label{def:grammar}
A completed expression is generated by
\begin{equation}\label{eq:grammar}
 E ::= hh\ \mid\ \Turn(E)\ \mid\ E_1E_2\ \mid\ hE\ \mid\ Eh.
\end{equation}
Every occurrence of $h$ denotes a distinct crossing circle and an independent auxiliary color. A closed expression in the main theorem is the cylindrical closure $C(FG)$, where both $F$ and $G$ are completed. Its ribbon framing, local disk framings and one global cut are retained throughout.
\end{definition}

The equivalent convention allowing every $h^s$, $s\ge2$, as an initial block gives the same grammar. Arbitrary finite nesting and repeated chains such as $E_{n+1}=h\Turn(E_n)$, $E_0=hh$, are allowed. The restriction on the two root children is substantive: strict feasibility alone does not give the selected high-index complex packet needed below.

At leaf $\ell$, choose $\eps_\ell\in\{0,1,-1\}$ and retain the corresponding physical half twist. Filling the crossing circle by the original disk-framed power $R^{-2p_\ell}$ replaces that leaf by the actual rational tangle
\begin{equation}\label{eq:rational-leaf}
 \Turn\bigl(R^{\eps_\ell-2p_\ell}\bigr).
\end{equation}
Choose the signs so that the endpoint pairing gives one main component and no closed component inside the tangle tree. An elementary pairing induction, proved in Section~\ref{sec:pairings}, supplies such a choice. After every auxiliary filling, no crossing-circle component remains. Products and quarter-turns of the rational tangles in \eqref{eq:rational-leaf} give a Conway algebraic link; with the chosen pairing it is a knot. This is a membership statement derived from actual tangle operations. Equality of source scalars will never be used as a substitute for tangle isotopy.

\begin{theorem}[Long-twist volume limit]\label{thm:main}
Fix a finite expression $C(FG)$ as in Definition~\ref{def:grammar} and a half-twist pattern $\eps$ giving a knot. There is a finite number $P(F,G,\eps)$ such that, for every independently chosen tuple of signed integers $\mathbf p=(p_\ell)$ with $|p_\ell|>P(F,G,\eps)$, the resulting knot $K_{F,G,\eps}(\mathbf p)$ is hyperbolic and
\begin{equation}\label{eq:main-volume}
 \lim_{N\to\infty}\frac{2\pi}{N}
 \log\left|\langle K_{F,G,\eps}(\mathbf p)\rangle_N\right|
 =\Vol\bigl(S^3\setminus K_{F,G,\eps}(\mathbf p)\bigr).
\end{equation}
The level tends to infinity through all integers, including even and composite levels. The expression, signs and twists are fixed before that limit. The threshold depends on the fixed expression and sign pattern, but not on a comparison between the magnitudes or signs of the individual sufficiently long twists.
\end{theorem}

The notation $\langle\cdot\rangle_N$ denotes the enhanced reduced invariant with unknot value one specified in Chapter~\ref{ch:source}. The full modified one-cut normalization differs by the single factor $\ii^{N+3}/N$ and a unit framing character, and therefore has the same limit. There is only one such modified factor, on the marked main component. The auxiliary circles use ordinary pivotal traces.

The already established isolated closures $C(h^s)$, $s\ge3$, can be included as a separate inherited family. The case $C(h^3)$ is described explicitly in Section~\ref{sec:isolated}. The isolated $C(hh)$ is excluded: its even-level source is zero, and its odd-level source is confined to a flat boundary. More general root-leaf closures $C(hE)$ are not included in Theorem~\ref{thm:main}. For example, in $C(h\Turn(h^s))$ the inner complete leaf indices equal $1/4$, at the boundary of the uniform packet used here.

\section{Proof architecture}
There are four logically distinct comparisons. First, finite matrix identities express the original scalar in a nonsemisimple polynomial algebra. Second, the complete positive source has a strict entropy maximizer, and explicit ideal cells identify its action with the volume of the actual augmented-link complement. Third, rigorous alias bounds and a finite complex lattice saddle produce a nonzero packet for the original mixed source. Finally, the exact filling transform selects a nonzero charge contribution whose critical action is the volume after filling.

\begin{figure}[ht]
\centering
\begin{tikzpicture}[node distance=8mm, every node/.style={font=\small}]
\node[box,text width=11.5cm] (s) {Evaluate the augmented tangle by its finite matrices\\Keep both polynomial values and derivatives at repeated roots};
\node[box,below=of s,text width=11.5cm] (p) {Maximize the positive complete-color sum\\Build its ideal cells and recover the actual augmented link metric};
\node[box,below=of p,text width=11.5cm] (a) {Bound the signed complementary color terms throughout the construction\\Obtain a nonzero holomorphic asymptotic near the complete maximum};
\node[box,below=of a,text width=11.5cm] (f) {Evaluate the entire oscillatory filling sum\\Keep contour boundary terms and identify the primitive filling slope};
\node[box,below=of f,text width=11.5cm,fill=green!5] (v) {The actual knot complement after all auxiliary fillings\\$N^{-1}\log|\langle K\rangle_N|\longrightarrow\Vol(S^3\setminus K)/(2\pi)$};
\draw[arr] (s)--(p);\draw[arr] (p)--(a);\draw[arr] (a)--(f);\draw[arr] (f)--(v);
\end{tikzpicture}

\caption{The proof keeps the finite source and the geometric marking together. A complete-color volume identity alone would not justify the final arrow.}\label{fig:architecture}
\end{figure}

Each comparison has an obstruction that explains its formulation. Repeated Casimir roots prevent a replacement by values at distinct nodes. The complete circle is nilpotent at even levels, so a fixed complete-color operator cannot be inverted to perform filling. Positive coefficients control the complete action but not the signed aliases. An asymptotic expansion at a single real color does not control a complex saddle displaced by a fixed long filling. Finally, the full filling period has both endpoint currents and infinitely many Fourier charges. The proof retains these features instead of hiding them in a normalization or a formal stationary-phase argument.

\section{Conventions and proof records}
All logarithms on a positive real interval are real. Their complex continuations are selected from that interval and are used only on explicitly chosen zero-free neighborhoods. Put $\arccot X\in(0,\pi)$ for real $X$, and write $\ii^2=-1$. A constant may depend on the fixed tangle tree, its chosen strict stationary point and a fixed compact analytic neighborhood, but never on $N$. A polynomial factor means $N^C$ for such a fixed constant $C$.

The appendices give the finite identities and the analytic estimates used in the main proof. The exact rational certificates are finite mathematical objects, rather than fitted numerical data: their verifier checks coverage, supporting inequalities and every outward-rounding remainder. Chapter~\ref{ch:reproduction} separates these checks from the continuum arguments. Internal proof review and successful execution are reported as such; neither is described as external peer review.
